% Job posting: https://ca.linkedin.com/jobs/view/embedded-software-designer-%E2%80%93-processing-algorithms-at-kepler-communications-inc-4454747734
\documentclass[11pt]{article}
\usepackage[left=1in,top=1in,right=1in,bottom=1in]{geometry}
\usepackage[hidelinks]{hyperref}
\usepackage{setspace}
\usepackage[T1]{fontenc}
\usepackage{newtxtext,newtxmath}
\ifdefined\pdfgentounicode
  \input{glyphtounicode}
  \pdfgentounicode=1
\fi
\setstretch{1.1}
\pagenumbering{gobble}
\begin{document}
\pagestyle{empty}

\noindent
Nishal Stanislaus Silva, PhD \\
Toronto, ON \\
nishal.silva@hotmail.com \,|\, +1 613 200 2030 \,|\, nishal.xyz \,|\, linkedin.com/in/nishal-silva

\vspace{1.5em}
\noindent Dear Hiring Manager,

\vspace{1em}
\noindent I am writing to apply for the Embedded Software Designer -- Processing \& Algorithms role at Kepler Communications. My graduate work and postdoctoral research have centered on exactly this problem: designing signal-processing algorithms that extract reliable information from noisy, continuous, real-world signals and run them under hard real-time constraints on genuinely constrained hardware. My MSc thesis and resulting IEEE Asilomar 2018 paper applied the S-transform to time-frequency signal extraction for onset detection in noisy audio streams. As a postdoctoral researcher at the University of Trento, I designed and shipped a real-time pattern-detection algorithm running at 14 ms inference on a Raspberry Pi 4, F1 0.76 at 31.4\% CPU utilization -- 74x faster than the 1038 ms DTW baseline it replaced -- published in the Journal of the Audio Engineering Society (2026).

\vspace{1em}
\noindent Beyond algorithm design, I have direct experience engineering under the kind of hardware constraints Kepler's satellite platforms impose. At McGill University's IDMIL/CIRMMT lab, I built and profiled a self-powered, battery-constrained embedded system, validating real-time feasibility against a strict compute and power budget -- an exercise directly analogous to designing algorithms for spacecraft power envelopes. I also built MIRaaS, a UDP-based edge-to-server offload architecture with characterized end-to-end latency, accepted at IEEE IS2 2026, which speaks to networked real-time signal transmission under bandwidth and latency limits. My primary engineering language is C++17, and my hardware-integration experience spans IMU/I2C sensor fusion, ADC/DAC audio hardware, and PLC interfacing on production manufacturing lines at MAS Holdings, where I also introduced CI/CD and code review practices to an embedded engineering team.

\vspace{1em}
\noindent I have not worked in the aerospace or satellite-communications industry directly, and I have no radiation-hardened or space-qualified hardware experience. What transfers directly, however, is the discipline: designing, benchmarking, and shipping signal-processing algorithms against hard latency and power budgets on constrained embedded platforms, and validating those designs with rigorous, reproducible measurement rather than assumption.

\vspace{1em}
\noindent I am a Canadian Permanent Resident, eligible to work in Canada without sponsorship, based in Toronto and available immediately. I would welcome the opportunity to discuss how my embedded signal-processing background can contribute to Kepler's processing and algorithms team.

\vspace{1.5em}
\noindent Sincerely, \\
Nishal Stanislaus Silva

\end{document}
